\subsection{Overview}

\note{This section is unfinished, but it wishes to introduce the following items: In Section 5.2, we will show that using the permutation representation, we may enumerate induced subgraphs of the comparability grid. This allows us to answer questions like ``Is the 3 by 3 Grid every an induced subgraph of some comp grid?'' Then, we will show how to, given a graph, determine if it is an induced subgraph of the comparability grid. The way we do this is by encoding the graph's data as a linear program, and checking if the LP is satisfiable. We conclude the section with a discussion to determine if known subgraph of the comp grid is tileable. This can be done by tiling the subgraph, and then using the LP code to check if it is still a subgraph. The LP code is constructive, so it will also output the construction if this is true. We remark that determining if a given graph is a vertex minor of some cgrid is an extremely expensive operation - it is NP-complete to check if its the vertex minor of a specific cgrid.}

\note{In section 5.3, we will introduce the idea of gadget extraction. Namely, give a subgraph that we know is tileable, how can we extract the circuit gadget we implement. We start with the simple example of a 3 wide grid with one additional edge every column, using purely circuit representation. We then show its ZX derivation, and say that because of the theorems in Section 2, we can find the gadget of any tileable subgraph efficiently. Now, given a Clifford circuit, can we determine if there is a graph state that implements it? I need a bit more time, but this may be possible using ZX.}

\note{In section 5.4, we will show the Lie algebra formalism for determining if a k-qubit gadget is k-qubit universal. There are some new developments in this area that I am excited to share, but this is mostly a copy-paste of the Lie Algebra notes Professor Grier made.}

\note{We then transition to the next section, where we ``present the fruits of our labor,'' or really in this case, the ``fruits of the labor of the programs.''}